% Part: lambda-calculus
% Chapter: introduction
% Section: basic-pr-lambda

\documentclass[../../../include/open-logic-section]{subfiles}

\begin{document}

\olfileid{lam}{int}{bas}
\olsection{Fungsi Rekursif Primitif Dhasar iku \usetoken{S}{lambda definable}}

\begin{lem}
Fungsi $\Zero$, $\Succ$, lan $\Proj{n}{i}$
iku !!{lambda definable}.
\end{lem}

\begin{proof}
% OLPL-234: Fungsi nol iku siji argumen; argumen u ditampa nanging ora digunakake.
$\Zero$ iku suku
$\lambd[u][\lambd[x][\lambd[y][y]]]$:
kanggo saben input, asile numeral Church nol.

Fungsi panerus~$\Succ$ ditetepake kanthi
$\fn{Succ}(u) = \lambd[x][\lambd[y][x(uxy)]]$.
Pikirna apa sebabe iki bener; kanggo saben
numeral $\num{n}$ kang dianggep iterator,
lan saben fungsi~$f$,
$\fn{Succ}(\num{n},f)$ iku fungsi kang,
yen diwenehi input~$y$, ngetrapake $f$
ping $n$ diwiwiti saka~$y$, banjur
ngetrapake maneh sepisan.

Ora perlu andharan dawa bab proyeksi:
$\fn{Proj}^n_i(x_0, \dots,
x_{n-1}) = x_i$. Tegese, miturut konvensi kita,
$\fn{Proj}^n_i$ iku suku lambda
$\lambd[x_0][\dots \lambd[x_{n-1}][x_i]]$.
\end{proof}

\end{document}
